\section{Results and Discussion}
\label{sec:results}

\subsection{Primary architecture comparison}

Table~\ref{tab:architecture_winners} identifies the lower-MAE architecture
at every cumulative checkpoint. All values are equal-site-weighted test
macro-MAEs, computed from the same origins and targets within each task.
Daily and monthly values evaluate the five-seed forecast ensemble; hourly
values evaluate the registered seed 42. A cumulative horizon includes all
leads up to that checkpoint, rather than only the final lead.

At the terminal horizons, TimesNet led the 72-hour task, with
52.48 versus 52.96~W/m$^2$ for DilatedRNN. DilatedRNN led the 30-day task
(43.85 versus 44.08~W/m$^2$), the principal one-month task
(13.01 versus 13.89~W/m$^2$), and the exploratory six-month task
(14.37 versus 14.46~W/m$^2$). Thus, the terminal winner was TimesNet for
the hourly contract and DilatedRNN for the daily and monthly contracts.
The largest relative terminal reduction was the one-month DilatedRNN gain:
6.38\% relative to TimesNet, calculated from unrounded MAEs. Counting these
four wins does not define an overall score because the targets, forecast
origins, and aggregation scales differ across tasks.

\begin{table*}[!t]
\centering
\caption{Focused architecture comparison: cumulative test macro-MAE in W/m$^2$. Bold marks the lower unrounded MAE within each row. $\Delta=\mathrm{MAE}_{\mathrm{TimesNet}}-\mathrm{MAE}_{\mathrm{DilatedRNN}}$ is computed before rounding. Hourly estimates use seed 42; daily and monthly estimates use five-seed ensembles. The complete reference-model benchmark is retained in Supplementary Table S1.}
\label{tab:architecture_winners}
\small
\setlength{\tabcolsep}{5pt}
\begin{tabular}{llrrrl}
\toprule
Task & Cumulative horizon & TimesNet & DilatedRNN & $\Delta$ & Lower MAE \\
\midrule
Hourly & 24 h & \textbf{49.39} & 50.73 & -1.35 & TimesNet \\
Hourly & 48 h & 52.16 & \textbf{52.15} & +0.02 & DilatedRNN \\
Hourly & 72 h & \textbf{52.48} & 52.96 & -0.48 & TimesNet \\
\addlinespace
Daily & 1 d & \textbf{43.20} & 43.27 & -0.07 & TimesNet \\
Daily & 3 d & \textbf{43.63} & 43.80 & -0.17 & TimesNet \\
Daily & 7 d & \textbf{43.95} & 44.08 & -0.13 & TimesNet \\
Daily & 14 d & \textbf{44.09} & 44.11 & -0.02 & TimesNet \\
Daily & 30 d & 44.08 & \textbf{43.85} & +0.23 & DilatedRNN \\
\addlinespace
Monthly 12--1 & 1 mo & 13.89 & \textbf{13.01} & +0.89 & DilatedRNN \\
\addlinespace
Monthly 12--6 & 1 mo & \textbf{14.86} & 14.89 & -0.03 & TimesNet \\
Monthly 12--6 & 3 mo & 14.49 & \textbf{14.33} & +0.16 & DilatedRNN \\
Monthly 12--6 & 6 mo & 14.46 & \textbf{14.37} & +0.09 & DilatedRNN \\
\bottomrule
\end{tabular}
\end{table*}


Table~\ref{tab:metricas-terminais} complements MAE with RMSE, nRMSE, and
$R^2$. The terminal RMSE ordering agrees with the MAE ordering: at 72 hours,
TimesNet obtained 101.17~W/m$^2$ versus 104.93~W/m$^2$ for DilatedRNN;
DilatedRNN had the lower RMSE in the daily and monthly tasks. These are
matched-test results, and their role is to resolve the architecture
comparison under the specified protocol. Results against all other
evaluated models are retained in the Supplementary Material and discussed
in Section~\ref{sec:reference_benchmark}.

\begin{table*}[htbp]
\centering
\caption{Terminal cumulative metrics for TimesNet and DilatedRNN. MAE and RMSE are in W/m$^2$; nRMSE is reported as a percentage.}
\label{tab:metricas-terminais}
\small
\setlength{\tabcolsep}{5pt}
\begin{tabular}{ll l rrrr}
\toprule
Task & Horizon & Model & MAE & RMSE & nRMSE (\%) & $R^2$ \\
\midrule
Hourly 336--72 & 72 h & TimesNet & 52.48 & 101.17 & 50.41 & 0.87 \\
Hourly 336--72 & 72 h & DilatedRNN & 52.96 & 104.93 & 52.31 & 0.86 \\
\addlinespace
Daily 365--30 & 30 d & TimesNet & 44.08 & 55.16 & 26.55 & 0.53 \\
Daily 365--30 & 30 d & DilatedRNN & 43.85 & 54.79 & 26.38 & 0.54 \\
\addlinespace
Monthly 12--1 & 1 mo & TimesNet & 13.89 & 17.49 & 8.51 & 0.88 \\
Monthly 12--1 & 1 mo & DilatedRNN & 13.01 & 16.02 & 7.86 & 0.91 \\
\addlinespace
Monthly 12--6 & 6 mo & TimesNet & 14.46 & 18.24 & 7.77 & 0.85 \\
Monthly 12--6 & 6 mo & DilatedRNN & 14.37 & 17.47 & 7.47 & 0.85 \\
\bottomrule
\end{tabular}
\end{table*}


\subsection{Site-paired differences and horizon dependence}

Table~\ref{tab:vitorias-timesnet-dilatedrnn} gives the site-paired comparison;
a negative difference denotes lower TimesNet MAE. TimesNet led at 24 and
72 hours by 1.35 and 0.48~W/m$^2$, respectively; DilatedRNN led at 48 hours
by 0.02~W/m$^2$. Differences are calculated before rounding individual
MAEs. TimesNet won seven sites at 24 hours and four at both 48 and 72 hours.
Its 72-hour macro advantage therefore reflects the magnitude of local gains,
even though DilatedRNN won more sites.

In the daily task, TimesNet led by 0.07--0.17~W/m$^2$ through seven days and
by 0.02~W/m$^2$ at 14 days. At 30 days, DilatedRNN led by
0.23~W/m$^2$ and won six sites. DilatedRNN also led the principal one-month
task by 0.89~W/m$^2$ and won six sites. In the exploratory six-month task,
TimesNet led the one-month prefix by 0.03~W/m$^2$ despite winning only four
sites. DilatedRNN led by 0.16 and 0.09~W/m$^2$ at three and six months,
respectively, and won seven sites at the terminal horizon. At three months,
TimesNet won six sites despite its higher macro-MAE. These reversals show
why both the magnitude and the distribution of site-level errors matter.

\begin{table*}[htbp]
\centering
\caption{Descriptive site-paired cumulative comparison. $\Delta=\mathrm{MAE}_{\mathrm{TimesNet}}-\mathrm{MAE}_{\mathrm{DilatedRNN}}$ in W/m$^2$; negative values favor TimesNet (TN), positive values favor DilatedRNN (DR). Wins count sites, not forecast origins. Overlapping multi-step targets are not independent observations.}
\label{tab:vitorias-timesnet-dilatedrnn}
\small
\setlength{\tabcolsep}{5pt}
\begin{tabular*}{\textwidth}{@{\extracolsep{\fill}}llrrrrrr@{}}
\toprule
Task & Horizon & Sites & Mean $\Delta$ & Median $\Delta$ & TN wins & DR wins & Ties \\
\midrule
Hourly 336--72 & 24 h & 10 & -1.35 & -1.33 & 7 & 3 & 0 \\
Hourly 336--72 & 48 h & 10 & 0.02 & 0.62 & 4 & 6 & 0 \\
Hourly 336--72 & 72 h & 10 & -0.48 & 0.50 & 4 & 6 & 0 \\
\addlinespace
Daily 365--30 & 1 d & 10 & -0.07 & -0.04 & 5 & 5 & 0 \\
Daily 365--30 & 3 d & 10 & -0.17 & -0.15 & 5 & 5 & 0 \\
Daily 365--30 & 7 d & 10 & -0.13 & -0.23 & 6 & 4 & 0 \\
Daily 365--30 & 14 d & 10 & -0.02 & -0.09 & 6 & 4 & 0 \\
Daily 365--30 & 30 d & 10 & 0.23 & 0.09 & 4 & 6 & 0 \\
\addlinespace
Monthly 12--1 & 1 mo & 10 & 0.89 & 0.35 & 4 & 6 & 0 \\
\addlinespace
Monthly 12--6 & 1 mo & 10 & -0.03 & 0.38 & 4 & 6 & 0 \\
Monthly 12--6 & 3 mo & 10 & 0.16 & -0.08 & 6 & 4 & 0 \\
Monthly 12--6 & 6 mo & 10 & 0.09 & 0.50 & 3 & 7 & 0 \\
\bottomrule
\end{tabular*}
\end{table*}


Figure~\ref{fig:local_heterogeneity} retains all ten terminal site contrasts.
The mixed signs show geographic heterogeneity and identify the locations at
which the task-level winner loses. They do not support a universal
architecture ranking.

\begin{figure*}[!t]
\centering
\includegraphics[width=0.90\textwidth]{heterogeneidade_local_delta_mae.png}
\caption{Site-level terminal cumulative contrast
$\Delta\mathrm{MAE}=\mathrm{MAE}_{\mathrm{TimesNet}}-
\mathrm{MAE}_{\mathrm{DilatedRNN}}$ in W/m$^2$. Negative (blue) cells
favor TimesNet; positive (red) cells favor DilatedRNN. Columns correspond to
72~h, 30~d, 1~month, and 6~months. Comparisons are descriptive; overlapping
multi-step target windows are not treated as independent observations.}
\label{fig:local_heterogeneity}
\end{figure*}

\subsection{Variability across seeds}

Daily and monthly learned models use five registered random seeds.
The 30-day seed-level macro-MAE was $44.29\pm0.29$~W/m$^2$ for TimesNet
and $44.10\pm0.60$~W/m$^2$ for DilatedRNN. At one month it was
$14.62\pm0.74$ versus $13.45\pm0.35$~W/m$^2$, respectively; at six months,
$15.24\pm0.76$ versus $14.77\pm0.45$~W/m$^2$.
Figure~\ref{fig:seed_variability} reports means and sample standard
deviations. These statistics need not equal
Table~\ref{tab:architecture_winners}: the latter evaluates the ensemble
forecast, whereas the former averages separately evaluated seed forecasts.
The hourly estimate uses one seed and cannot quantify seed variability.

The 0.23~W/m$^2$ terminal daily ensemble gap is smaller than either
architecture's seed-level standard deviation; the six-month gap is only
0.09~W/m$^2$. Comparing runs with the same seed, DilatedRNN had lower
terminal MAE in three of five daily runs and four of five runs in each
monthly task. The one-month advantage is thus more consistent across the
retained initializations, but is not universal. Seed standard deviations
measure training variability on the same test year; they are neither
confidence intervals nor estimates of uncertainty across future years.
Comparing an ensemble gap with these deviations is descriptive, not a
significance test.

\begin{figure*}[!t]
\centering
\includegraphics[width=\textwidth]{variabilidade_arquiteturas.png}
\caption{TimesNet and DilatedRNN seed-level macro-MAE: means with error bars
spanning one sample standard deviation. Color and marker shape identify
each architecture. The hourly extension is shown as $n=1$ without an error
bar. The full learned-model comparison is retained in the Supplementary Material.}
\label{fig:seed_variability}
\end{figure*}

\subsection{Reference benchmark and relation to the source protocols}
\label{sec:reference_benchmark}

The full benchmark remains necessary to interpret the focused comparison.
Supplementary Table S1 and Figure S1 retain every evaluated reference model
and cumulative checkpoint. XGBoost had the lowest macro-MAE at all hourly
checkpoints (50.50~W/m$^2$ at 72 hours) and at 1, 3, and 14 days.
LSTM led at 7 days (43.69~W/m$^2$), and DilatedRNN led at 30 days
(43.85~W/m$^2$), ahead of XGBoost (43.93~W/m$^2$).
Calendar climatology led both monthly tasks, with 12.57~W/m$^2$ at one month
and 13.03~W/m$^2$ over six months. Hourly climatology was not evaluated;
its missing entries are not zero errors. Consequently, a win over the other
focal architecture must not be read as a win over the entire benchmark.

The source studies motivate the hourly TimesNet and monthly DilatedRNN
contracts, but their original error values are not a matched comparison.
The present evaluation aligns inputs, origins, post-processing, aggregation,
and metrics within each task. Its ranking therefore uses the matched
predictions throughout, rather than combining results from the two
source protocols.

\subsection{Favorable and unfavorable cases}

The recorded global case rule selects the maximum and minimum, over all
available sites, origins, and declared cumulative horizons, of
$\mathrm{MAE}_{\mathrm{DilatedRNN}}-\mathrm{MAE}_{\mathrm{TimesNet}}$.
Figure~\ref{fig:contrasting_cases} plots both selected forecasts.
Weather is attached after selecting the error extremes. The largest TimesNet
gain occurred at GM Factory Zero for the 24-hour forecast beginning on
31 March 2024 in fixed local time: TimesNet and DilatedRNN obtained
45.43 and 109.90~W/m$^2$, respectively, a 64.47~W/m$^2$ advantage.
NASA POWER provided post-hoc context of 99.64\% mean cloud amount and
0.80~mm daily precipitation.

The largest deficit occurred at Tesla Gigafactory Texas for the 24-hour
forecast beginning on 28 March 2024: 150.02 versus 73.92~W/m$^2$, a
76.10~W/m$^2$ TimesNet disadvantage. NASA POWER reported 0.06~mm daily
precipitation and 2.14\% cloud amount; neither adverse-day threshold was
exceeded. An absent flag does not establish the absence of subdaily weather
variability. Both cases use the single hourly seed. Moreover, a global
selection based on raw MAE contrasts across resolutions can favor tasks
with larger absolute error scales. These examples cannot establish a
cross-resolution ranking or replace the complete aggregate evaluation.

\begin{figure*}[!t]
\centering
\includegraphics[width=\textwidth]{previsoes_casos_contrastantes.png}
\caption{Reference GHI and forecasts for the globally selected largest
TimesNet gain and deficit. Panels share a GHI scale. MAEs and descriptive
weather records refer to the selected 24-hour windows; meteorology was
joined only after case selection.}
\label{fig:contrasting_cases}
\end{figure*}

\subsection{Geographic and climatic interpretation}

Figure~\ref{fig:geographic_context} locates the ten sites in their descriptive
K\"oppen--Geiger context. With only ten purposively selected locations,
apparent differences between zones cannot establish a causal climatic effect.
Sites within a country can occupy distinct regimes, while nearby sites may
share correlated weather.

As an error-independent contextual example, the recorded meteorological
rule selected Hyundai Metaplant Georgia on 5 August 2024, with
82.66~mm/day corrected precipitation and 99.47\% cloud amount. This case
was selected without consulting either model's errors. Its meteorological
record is retained in the Supplementary Material, and its surrounding
forecast profiles are shown next.

\subsection{Multi-day irradiance profiles}
\label{sec:multiday_profiles}

Figure~\ref{fig:seven_day_profiles} extends the two error-selected cases
and the independent weather case to seven calendar days each: the event
date, three preceding days, and three following days. This fixed extension
rule does not optimize the displayed week using model errors. Each panel
contains 168 unique hourly targets and seven successive midnight forecast
origins. Only leads 1--24 of each 72-hour model output are retained; the
context advances daily while the fitted parameters remain fixed. The figure
therefore shows seven daily forecasts, not one 168-hour forecast.

Across these seven-day displays, TimesNet and DilatedRNN respectively
obtained MAEs of 99.73 and 115.07~W/m$^2$ at GM Factory Zero,
91.90 and 89.19~W/m$^2$ at Tesla Gigafactory Texas, and
109.42 and 122.98~W/m$^2$ at Hyundai Metaplant Georgia. These local
window errors are distinct from the single-event errors and the annual
macro-MAEs. The additional days reveal underprediction of some high-GHI
peaks and overprediction of low-GHI daytime periods by both architectures.
At the independently selected Georgia event, both forecasts exceed much
of the reference daytime profile. These trajectories expose limitations
of the GHI-only models, while the annual matched evaluation remains the
basis for declaring task-level winners.

\begin{figure*}[!t]
\centering
\includegraphics[width=\textwidth]{perfis_sete_dias.png}
\caption{Seven-day hourly GHI profiles around the two error-selected cases
and the error-independent NASA POWER event. Each curve joins seven
24-hour forecast prefixes issued on successive local midnights; it is not
a single seven-day forecast. Shading marks the original selected day.
MAEs use all 168 displayed targets per site, including nighttime; all
forecasts use seed 42 and the registered hourly post-processing.
Dates follow the fixed local-time index, and all panels share a GHI scale.}
\label{fig:seven_day_profiles}
\end{figure*}
